% ============================================================================
% FILE: sections/08_appendix.tex
% PURPOSE: รายละเอียดเพิ่มเติมที่ไม่ได้ใส่ในเนื้อหาหลัก
% ============================================================================

\appendix

\section{รายละเอียดการคำนวณเชิงกล}
% การคำนวณฉบับเต็ม Shaft, Belt, Bearing, Deflection
% อ้างอิงไฟล์ PDF ที่มีอยู่แล้วใน Project

\section{รายละเอียดการคำนวณ Motion Profile}
% การคำนวณ S-Curve ฉบับเต็ม ทุกช่วง (Acceleration, Constant, Deceleration)

\section{Schematic ไฟฟ้าฉบับเต็ม}
% Schematic ขนาดเต็มอ่านได้ง่าย

\begin{figure}[H]
  \centering
%   \includegraphics[width=0.95\textwidth]{images/schematic_full.png}
  % [ใส่ภาพ Schematic ขนาดเต็ม]
  \caption{Electrical Schematic ฉบับเต็ม}
  \label{fig:schematic_full}
\end{figure}

\section{Code Snippet หลักของ STM32 Firmware}
% Code สำคัญ เช่น PID Loop, Encoder Read, S-Curve Generator
% ใช้ lstlisting environment

\begin{lstlisting}[language=C, caption={ตัวอย่าง PID Control Loop บน STM32}]
/* [ใส่ Code จริงของ PID Loop] */
\end{lstlisting}

\section{ผลการ Simulation Simscape ฉบับเต็ม}
% Plot และผลลัพธ์จาก Simscape ทุก Scenario

\begin{figure}[H]
  \centering
%   \includegraphics[width=0.85\textwidth]{images/simscape_full_result.png}
  % [ใส่ Plot ทั้งหมดจาก Simscape: Position, Velocity, Control Effort]
  \caption{ผล Simulation Simscape ฉบับเต็ม}
  \label{fig:simscape_full}
\end{figure}

\section{Log การทดสอบ Unit Test}
% Serial Log, Oscilloscope Screenshot จากการทดสอบต่างๆ

\begin{figure}[H]
  \centering
%   \includegraphics[width=0.80\textwidth]{images/serial_log_encoder.png}
  % [ใส่ภาพ Serial Monitor Log แสดงค่า Encoder ขณะทดสอบ]
  \caption{Serial Log ขณะทดสอบการอ่านค่า Encoder}
  \label{fig:serial_log}
\end{figure}